% Nouveau EGA V (cote 25): archive cote 25, site page 128 (PDF page 129).
% typescript. Draft compared: none.
% AI-transcribed by Claude Opus 5.5 (claude-opus-5-5, adaptive thinking, effort high), Message Batches API, batch ega25b (msgbatch_01A7n5YuXubUPNK1zsepi8RG), 2026-11-13; not yet checked by eye.
\scanpage{25/128}
\mspage{31}

par une section inversible de $\underline{O}_S$ , donc $\phi$ et \unsure{$\psi$} définissent le même point
de $\overset{\hand{$\vee$}}{P}$ à valeurs dans S. Comme tout point de $\overset{\hand{$\vee$}}{P}$ à valeurs dans S est défini,
localement sur S, par une section $\phi$ de $\underline{E}_S$ qui ne s'annule pas (cf Chap II)
a) s'ensuit.   Pour prouver b), on note le

\uline{Lemme}   Soit D un syst lin de diviseurs sur X non vide loc de type fini
sur k alg clos, paramétré par $\overset{\hand{$\vee$}}{P}$(E), et considérons le morphisme correspon-%
dant  f: X-Z $\longrightarrow$ P , où Z est l'ensemble des points fixes. Alors \xout{xx} si
r = rang$_k$E $>$\hand{0,} il existe  r+1 point\hand{s}\inserted[typed]{$x_i$}\inserted{$(1\leq i\leq r+1)$} de \xout{(Xx} X(k)-Z(k)  tels que les
f(x$_i$) forment une  ``base projective'' de P, i.e. tels que  pour toute
partie J de [1,r+1] ayant \hand{r} \xout{mi} éléments,  les  f(x$_i$),i$\in$\hand{J}, \xout{xxxx} ne sont pas
contenus dans un hyperplan de P.

On peut évidemment supposer Z=$\emptyset$ , \deleted{X intègre (en le remplaçant
par une\xout{Xxxxxx} composante irréductible T de X muni de la structure induite
réduite).} \xout{OxxaxxxxxapplicationxxinjectivexxxE} Comme \inserted[typed]{par \over{4}{\unsure{3|5}}.1.} f(X) n'est contenu dans
aucun hyperplan de P, on en conclut d'abord l'existence de r points
$x_i$\inserted{$(1\leq i\leq r)$} de X \xout{telsxque}\inserted[typed]{tels que les f($x_i$) soient} projectivement indépendants dans P, i.e. soient définies
par des formes linéairement indépendantes sur E.   Reste à prouver qu'il
existe un  $x_{r+1} = x$\deleted{\illegible{1}}   dans X(k) tel que  f(x) ne soit dans aucun des r
hyperplans \inserted{$H_i$} définis par les \deleted{\unsure{f($x_i$)}} systèmes de r-1 parmi les f(x$_i$). Or dans
le cas contraire, compte tenu du sorite Jacobson,  on aurait
f(X) $\subset \bigcup_{\hand{$i$}}$ H$_i$ , donc \xout{Xxseraitxlaxréunionxdesxxxx}\xout{xxx} si X$_0$  est une composante
irréductible de X, f(X$_0$) serait contenu dans un des H$_i$ , ce qui contredit
\over{\unsure{4}}{\unsure{3|5}}.1. - Ceci posé,  pour prouver b), on peut  évidemment supposer Z=$\emptyset$ ,
et utilisant \over{\unsure{4}}{\unsure{3|5}}.1.  on est ramené à prouver qu'un automorphisme\inserted{u} de $\overset{\hand{$\vee$}}{P}$ est
déterminé quand on connait \xout{xxx}\inserted[typed]{le} composé\inserted[typed]{de\xout{\unsure{s}} son contragrédient $\overset{\hand{$\vee$}}{u}$ dans P} avec f: X$\longrightarrow$ P , - et que l'asser-%
tion analogue est vraie après tout changement de base S$\rightarrow$Spec(k)
pour un  automorphisme u de $\overset{\hand{$\vee$}}{P}_S$  .  Or cela résulte aussitôt  du lemme

\mspage{128}
